\documentclass{article}
\usepackage[utf8]{inputenc}
\usepackage{amsthm}

\usepackage{amsmath}
\usepackage{amssymb}
\usepackage{enumitem}
\usepackage{ marvosym }
\usepackage{amscd}

\newtheorem*{defn}{Definition}

\title{Up and down}
\author{Adam W. Grant }
\date{January 2021}

\begin{document}

\maketitle

\section{Introduction}
After defining what it means to be an up-set or a down set, the authors of \cite{davey_priestley_2002} introduce the reader to the $\uparrow$ and $\downarrow$ operators.
\begin{defn}
Given an arbitrary subset $Q$ of $P$ and $x \in P$, we define
\begin{align*} 
\downarrow Q := \{y \in P : \exists(x : x \in Q : y \leq x)\} &\text{ and }  \uparrow Q := \{y \in P : \exists(x : x \in Q : y \geq x)\} \\ 
\downarrow x := \{y \in P : y \leq x\} &\text{ and }  \uparrow x := \{y \in P : y \geq x)\}
\end{align*}
These are read `down Q', etc.
\end{defn}
In this note we will derive the coreflexive relations associated to each operation.

\section{Down Q, relationally}
By inspection, we see that the characteristic predicate of $\downarrow.Q$ is:
\begin{equation}
    \exists(x : x \in Q : y \leq x)
\end{equation}
Let's calculate its PFT
\begin{itemize}
    \item $\exists(x : x \in Q : y \leq x)$
    \item[$\equiv$] $\{\text{Trading; Coreflexives}\}$
    \item[] $\exists(x : : y (\leq \cdot \Phi_{Q}) x)$
    \item[$\equiv$] $\{\text{$\wedge$ Idempotent; Trading}\}$
    \item[] $\exists(x : y (\leq \cdot \Phi_{Q}) x : y (\leq \cdot \Phi_{Q}) x)$
    \item[$\equiv$] $\{\text{Converse}\}$
    \item[] $\exists(x : y (\leq \cdot \Phi_{Q}) x : x (\leq \cdot \Phi_{Q})^{\circ} y)$
    \item[$\equiv$] $\{\text{Composition}\}$
    \item[] $y ((\leq \cdot \Phi_{Q}) \cdot (\leq \cdot \Phi_{Q})^{\circ}) y$
    \item[$\equiv$] $\{\text{Image}\}$
    \item[] $y (img(\leq \cdot \Phi_{Q})) y$
    \item[\Heart]
\end{itemize}
Characteristic predicates are sometimes denoted by the greek letter $\chi$, so let's say that:
\begin{equation*}
    \chi.y \,\equiv\, y (img(\leq \cdot \Phi_{Q})) y
\end{equation*}
Then we can use the fact\footnote{Explained in reference \cite{Oliveira2009}} that
\begin{equation*}
    z \,\Phi_{\chi} y \,\equiv\, z = y \wedge \chi. y
\end{equation*}
to find the coreflexive relation associated to (1):
\begin{itemize}
    \item $z \,\Phi_{\chi} y$
    \item[--] $\chi.y \; \equiv \; y(img(\leq \cdot \Phi_{Q})) y$
    \item[$\equiv$] $\{\text{Coreflexives; $\chi$}\}$
    \item[] $z = y \wedge y(img(\leq \cdot \Phi_{Q})) y$
    \item[$\equiv$] $\{\text{Substitution}\}$
    \item[] $z = y \wedge z(img(\leq \cdot \Phi_{Q})) y$
    \item[$\equiv$] $\{\text{Introduce $id$; Meet}\}$
    \item[] $z (id \cap img(\leq \cdot \Phi_{Q}) y$
    \item[$\equiv$] $\{\text{Range}\}$
    \item[] $z (\rho(\leq \cdot \Phi_{Q}) y$
    \item[\Heart]
\end{itemize}
And so we have that the coreflexive relation associated to the set $\downarrow.Q$ is
\begin{equation}
    \rho(\leq \cdot \Phi_{Q})
\end{equation}
Thanks to the duality between the range and domain operators, we can express (3) equivalently as
\begin{equation}
    \delta(\Phi_{Q} \cdot \geq)
\end{equation}
which is perhaps more intuitive.

\section{Up Q, relationally}
The PFT for $\uparrow$ the same (mutatis mutandis) as that already seen for $\downarrow$, so we'll just state the result.  The coreflexive associated to the set $\uparrow.Q$ is:
\begin{equation}
    \rho(\geq \cdot \Phi_{Q})
\end{equation}
or
\begin{equation}
    \delta(\Phi_{Q} \cdot \leq)
\end{equation}
Relation (5) is somewhat less intuitive than (4).

\section{What about $x$?}
We can use what we already have to find the relational counterparts for principle down- and up-sets.  For the former, we substitute $\{x\}$ for $Q$ in (3)
\begin{equation}
    \delta(\Phi_{\{x\}} \cdot \geq)
\end{equation}
As for $\uparrow.x$, we perform the same substitution to (7)
\begin{equation}
    \rho(\geq \cdot \Phi_{\{x\}})
\end{equation}
There's no need to define separate constructions for sets versus individuals.

\bibliographystyle{plain}
\bibliography{refs}

\end{document}
